\documentclass[10pt]{article}
\usepackage[top=0.55in, bottom=0.5in, left=0.65in, right=0.65in, headheight=14pt, footskip=18pt]{geometry}
\usepackage{amsmath,amssymb,amsfonts}
\usepackage{graphicx}
\usepackage{xcolor}
\usepackage{tikz}
\usepackage{fancyhdr}
\usepackage[vietnamese]{babel}

% Định nghĩa màu sắc đặc trưng của Stewart Calculus
\definecolor{stewartcyan}{RGB}{0, 130, 165}
\definecolor{stewartrule}{RGB}{0, 160, 205}
\definecolor{darkgray}{RGB}{60, 60, 60}

% Thiết lập ký hiệu giới hạn luôn có chỉ số ở dưới (như bản gốc)
\let\oldlim\lim
\renewcommand{\lim}{\oldlim\limits}

% Cấu hình header & footer
\pagestyle{fancy}
\fancyhf{}
\renewcommand{\headrulewidth}{0pt}
\fancyhead[R]{\small\sffamily\textbf{\color{darkgray}MỤC 2.2}\quad \textmd{Giới hạn của hàm số}\qquad\textbf{93}}

% Biểu tượng thiết bị vẽ đồ thị (Graphing Device Icon)
\newcommand{\graphdeviceicon}{%
  \begin{tikzpicture}[baseline=-0.5ex, scale=0.85]
    \draw[stewartcyan, thick, rounded corners=1.2pt] (0,0) rectangle (0.42cm, 0.32cm);
    \draw[stewartcyan, thick] (0.06cm, 0.08cm) to[out=25, in=205] (0.36cm, 0.24cm);
  \end{tikzpicture}%
}

\begin{document}
\small
\setlength{\abovedisplayskip}{3pt}
\setlength{\belowdisplayskip}{3pt}

\noindent
\begin{minipage}[t]{0.482\textwidth}
  \noindent
  \textbf{(f)} Phương trình các đường tiệm cận đứng

  \vspace{0.3em}
  \noindent
  \centering
  \includegraphics[width=0.88\linewidth]{images/p0128_fig1.png}

  \vspace{0.5em}
  \raggedright
  \noindent
  \textbf{10.} Một bệnh nhân được tiêm một liều thuốc 150-mg sau mỗi 4 giờ. Đồ thị cho thấy lượng thuốc $f(t)$ trong máu sau $t$ giờ. Hãy tìm
  \[
    \lim_{t \to 12^-} f(t) \qquad\text{và}\qquad \lim_{t \to 12^+} f(t)
  \]
  và giải thích ý nghĩa thực tế của các giới hạn một bên này.

  \vspace{0.4em}
  \noindent
  \centering
  \includegraphics[width=0.72\linewidth]{images/p0128_fig2.png}

  \vspace{0.6em}
  \raggedright
  \noindent
  \textbf{\textsf{\color{stewartcyan}11--12}}\ Phác thảo đồ thị của hàm số và sử dụng nó để xác định các giá trị của $a$ mà tại đó $\lim_{x \to a} f(x)$ tồn tại.

  \vspace{0.4em}
  \noindent
  \textbf{11.} $f(x) = \begin{cases} 
    e^x & \text{nếu } x \le 0 \\ 
    x - 1 & \text{nếu } 0 < x < 1 \\ 
    \ln x & \text{nếu } x \ge 1 
  \end{cases}$

  \vspace{0.5em}
  \noindent
  \textbf{12.} $f(x) = \begin{cases} 
    \sqrt[3]{x} & \text{nếu } x \le -1 \\ 
    x & \text{nếu } -1 < x \le 2 \\ 
    (x - 1)^2 & \text{nếu } x > 2 
  \end{cases}$

  \vspace{0.7em}
  \noindent
  \graphdeviceicon\;\textbf{\textsf{\color{stewartcyan}13--14}}\ Sử dụng đồ thị của hàm số $f$ để xác định giá trị của mỗi giới hạn sau, nếu nó tồn tại. Nếu không tồn tại, hãy giải thích lý do.
  
  \vspace{0.3em}
  \noindent
  (a) $\lim_{x \to 0^-} f(x)$ \qquad (b) $\lim_{x \to 0^+} f(x)$ \qquad (c) $\lim_{x \to 0} f(x)$

  \vspace{0.4em}
  \noindent
  \begin{tabular}{@{}p{0.46\linewidth}@{\hspace{0.04\linewidth}}p{0.50\linewidth}@{}}
    \textbf{13.} $f(x) = x\sqrt{1 + x^{-2}}$ & \textbf{14.} $f(x) = \dfrac{e^{1/x} - 2}{e^{1/x} + 1}$
  \end{tabular}

  \vspace{0.7em}
  \noindent
  \textbf{\textsf{\color{stewartcyan}15--18}}\ Phác thảo đồ thị của một ví dụ về hàm số $f$ thỏa mãn tất cả các điều kiện đã cho.

  \vspace{0.4em}
  \noindent
  \textbf{15.} $\lim_{x \to 1^-} f(x) = 3, \quad \lim_{x \to 1^+} f(x) = 0, \quad f(1) = 2$

  \vspace{0.45em}
  \noindent
  \textbf{16.} $\lim_{x \to 0} f(x) = 4, \quad \lim_{x \to 8^-} f(x) = 1, \quad \lim_{x \to 8^+} f(x) = -3,$ \\[0.25em]
  \phantom{\textbf{16. }}$f(0) = 6, \quad f(8) = -1$
\end{minipage}%
\hfill
\begin{minipage}[t]{0.482\textwidth}
  \noindent
  \textbf{17.} $\lim_{x \to -1^-} f(x) = 0, \quad \lim_{x \to -1^+} f(x) = 1, \quad \lim_{x \to 2} f(x) = 3,$ \\[0.25em]
  \phantom{\textbf{17. }}$f(-1) = 2, \quad f(2) = 1$

  \vspace{0.45em}
  \noindent
  \textbf{18.} $\lim_{x \to -3^-} f(x) = 3, \quad \lim_{x \to -3^+} f(x) = 2, \quad \lim_{x \to 3^-} f(x) = -1,$ \\[0.25em]
  \phantom{\textbf{18. }}$\lim_{x \to 3^+} f(x) = 2, \quad f(-3) = 2, \quad f(3) = 0$

  \vspace{0.4em}
  \noindent{\color{stewartrule}\rule{\linewidth}{0.55pt}}
  \vspace{0.25em}

  \noindent
  \textbf{\textsf{\color{stewartcyan}19--22}}\ Dự đoán giá trị của giới hạn (nếu tồn tại) bằng cách tính giá trị của hàm số tại các số đã cho (chính xác đến sáu chữ số thập phân).

  \vspace{0.35em}
  \noindent
  \textbf{19.} $\lim_{x \to 3} \dfrac{x^2 - 3x}{x^2 - 9}$, \\[0.2em]
  \phantom{\textbf{19. }}$x = 3.1, 3.05, 3.01, 3.001, 3.0001,$ \\[0.15em]
  \phantom{\textbf{19. }}$2.9, 2.95, 2.99, 2.999, 2.9999$

  \vspace{0.35em}
  \noindent
  \textbf{20.} $\lim_{x \to -3} \dfrac{x^2 - 3x}{x^2 - 9}$, \\[0.2em]
  \phantom{\textbf{20. }}$x = -2.5, -2.9, -2.95, -2.99, -2.999, -2.9999,$ \\[0.15em]
  \phantom{\textbf{20. }}$-3.5, -3.1, -3.05, -3.01, -3.001, -3.0001$

  \vspace{0.35em}
  \noindent
  \textbf{21.} $\lim_{t \to 0} \dfrac{e^{5t} - 1}{t}, \quad t = \pm 0.5, \pm 0.1, \pm 0.01, \pm 0.001, \pm 0.0001$

  \vspace{0.35em}
  \noindent
  \textbf{22.} $\lim_{h \to 0} \dfrac{(2 + h)^5 - 32}{h}$, \\[0.2em]
  \phantom{\textbf{22. }}$h = \pm 0.5, \pm 0.1, \pm 0.01, \pm 0.001, \pm 0.0001$

  \vspace{0.4em}
  \noindent{\color{stewartrule}\rule{\linewidth}{0.55pt}}
  \vspace{0.25em}

  \noindent
  \textbf{\textsf{\color{stewartcyan}23--28}}\ Sử dụng một bảng giá trị để ước tính giá trị của giới hạn. Nếu bạn có một thiết bị vẽ đồ thị, hãy sử dụng nó để xác nhận kết quả bằng hình ảnh.

  \vspace{0.35em}
  \noindent
  \begin{tabular}{@{}p{0.48\linewidth}@{\hspace{0.04\linewidth}}p{0.48\linewidth}@{}}
    \textbf{23.} $\lim_{x \to 4} \dfrac{\ln x - \ln 4}{x - 4}$ & \textbf{24.} $\lim_{p \to -1} \dfrac{1 + p^9}{1 + p^{15}}$ \\[0.65em]
    \textbf{25.} $\lim_{\theta \to 0} \dfrac{\sin 3\theta}{\tan 2\theta}$ & \textbf{26.} $\lim_{t \to 0} \dfrac{5^t - 1}{t}$ \\[0.65em]
    \textbf{27.} $\lim_{x \to 0^+} x^x$ & \textbf{28.} $\lim_{x \to 0^+} x^2 \ln x$
  \end{tabular}

  \vspace{0.4em}
  \noindent{\color{stewartrule}\rule{\linewidth}{0.55pt}}
  \vspace{0.25em}

  \noindent
  \textbf{\textsf{\color{stewartcyan}29--40}}\ Xác định giới hạn vô cực.

  \vspace{0.35em}
  \noindent
  \begin{tabular}{@{}p{0.48\linewidth}@{\hspace{0.04\linewidth}}p{0.48\linewidth}@{}}
    \textbf{29.} $\lim_{x \to 5^+} \dfrac{x + 1}{x - 5}$ & \textbf{30.} $\lim_{x \to 5^-} \dfrac{x + 1}{x - 5}$ \\[0.7em]
    \textbf{31.} $\lim_{x \to 2} \dfrac{x^2}{(x - 2)^2}$ & \textbf{32.} $\lim_{x \to 3^-} \dfrac{\sqrt{x}}{(x - 3)^5}$ \\[0.7em]
    \textbf{33.} $\lim_{x \to 1^+} \ln(\sqrt{x} - 1)$ & \textbf{34.} $\lim_{x \to 0^+} \ln(\sin x)$ \\[0.7em]
    \textbf{35.} $\lim_{x \to (\pi/2)^+} \dfrac{1}{x} \sec x$ & \textbf{36.} $\lim_{x \to \pi^-} x \cot x$ \\[0.7em]
    \textbf{37.} $\lim_{x \to 1} \dfrac{x^2 + 2x}{x^2 - 2x + 1}$ & \textbf{38.} $\lim_{x \to 3^-} \dfrac{x^2 + 4x}{x^2 - 2x - 3}$ \\[0.7em]
    \textbf{39.} $\lim_{x \to 0} (\ln x^2 - x^{-2})$ & \textbf{40.} $\lim_{x \to 0^+} \left(\dfrac{1}{x} - \ln x\right)$
  \end{tabular}
\end{minipage}

\vfill
\noindent
\begin{minipage}{\textwidth}
  \centering
  \fontsize{6.3pt}{7.5pt}\selectfont\color{gray}
  Chung bản quyền © 2021 Cengage Learning. Đã đăng ký mọi quyền. Không được phép sao chép, quét hoặc nhân bản toàn bộ hoặc một phần. WCN 02-200-203\\ 
  Bản quyền năm 2021 Cengage Learning. Tất cả các quyền được bảo lưu. Do quyền điện tử, một số nội dung bên thứ ba có thể bị ẩn khỏi sách điện tử và/hoặc chương điện tử.\\ 
  Đánh giá biên tập cho thấy nội dung bị ẩn không ảnh hưởng đáng kể đến trải nghiệm học tập tổng thể. Cengage Learning có quyền gỡ bỏ nội dung bổ sung vào bất kỳ lúc nào nếu các hạn chế quyền yêu cầu.
\end{minipage}

\end{document}